% ===== PRÉAMBULE CANONIQUE — à coller VERBATIM en tête de CHAQUE .tex =====
\documentclass[11pt,a4paper]{article}
\usepackage[utf8]{inputenc}\usepackage[T1]{fontenc}\usepackage{lmodern}
\usepackage[french]{babel}
\usepackage{amsmath,amssymb,mathtools}\usepackage{icomma}
\usepackage{siunitx}\sisetup{locale=FR}  % \qty{}{}, \si{}, \ang{}
\usepackage{xcolor}\usepackage{colortbl}
\usepackage{tikz}\usepackage{pgfplots}\pgfplotsset{compat=1.18}
\usepackage[siunitx,europeanresistors]{circuitikz} % circuits (élec.)
\usepackage[most]{tcolorbox}\usepackage{enumitem}\usepackage{array,booktabs}
\usepackage[a4paper,margin=2.2cm]{geometry}
\usetikzlibrary{arrows.meta,calc,angles,quotes,positioning,patterns,%
  backgrounds,shapes.geometric,decorations.pathmorphing,decorations.markings,intersections}
\definecolor{primary}{RGB}{222,49,99}\definecolor{primarydark}{RGB}{150,22,55}
\definecolor{defcol}{RGB}{0,114,178}\definecolor{methcol}{RGB}{52,92,148}
\definecolor{excol}{RGB}{0,135,90}\definecolor{attncol}{RGB}{200,40,55}
\tcbset{boxbase/.style={enhanced,breakable,boxrule=0.9pt,arc=2.5pt,
  left=3mm,right=3mm,top=2.3mm,bottom=2.3mm,fonttitle=\bfseries,coltitle=white}}
% --- boîtes TUTORIEL (noms EXACTS, ne pas renommer) ---
\newtcolorbox{cle}[1][]{boxbase,colback=defcol!6,colframe=defcol,
  title={\textsc{À retenir}\ifx\relax#1\relax\relax\else\textmd{ : \itshape #1}\fi}}
\newtcolorbox{methode}[1][]{boxbase,colback=methcol!7,colframe=methcol,sharp corners,
  title={\textsc{Méthode}\ifx\relax#1\relax\relax\else\textmd{ --- \itshape #1}\fi}}
\newtcolorbox{exemplebox}[1][]{boxbase,colback=excol!5,colframe=excol,
  title={\textsc{Exemple}\ifx\relax#1\relax\relax\else\textmd{ : \itshape #1}\fi}}
\newtcolorbox{piege}[1][]{boxbase,colback=attncol!6,colframe=attncol,
  title={\textsc{Piège}\ifx\relax#1\relax\relax\else\textmd{ : \itshape #1}\fi}}
\newtcolorbox{remarque}[1][]{boxbase,colback=black!4,colframe=black!45,
  title={\textsc{Remarque}\ifx\relax#1\relax\relax\else\textmd{ : \itshape #1}\fi}}
% --- boîtes EXERCICE / CORRIGÉ (noms EXACTS, auto-comptées) ---
\newtcolorbox[auto counter]{exo}[1][]{boxbase,colback=primary!4,colframe=primary,
  title={\textsc{Exercice}~\thetcbcounter\ifx\relax#1\relax\relax\else\textmd{ --- \itshape #1}\fi}}
\newtcolorbox[auto counter]{corrige}[1][]{boxbase,colback=excol!4,colframe=excol,
  title={\textsc{Corrigé}~\thetcbcounter\ifx\relax#1\relax\relax\else\textmd{ --- \itshape #1}\fi}}
\newcommand{\vv}[1]{\vec{#1}}\newcommand{\ds}{\displaystyle}
\newcommand{\R}{\mathbb{R}}\newcommand{\N}{\mathbb{N}}\newcommand{\Z}{\mathbb{Z}}
% =========================================================================
\begin{document}
\usetikzlibrary{babel}
\begin{exo}[loi d'Ohm : trouver la tension]
Une résistance $R=\qty{20}{\ohm}$ est traversée par un courant $I=\qty{0.3}{\ampere}$. Calcule la
tension $U$ à ses bornes.
\end{exo}

\begin{exo}[loi d'Ohm : trouver l'intensité]
On applique une tension $U=\qty{12}{\volt}$ aux bornes d'une résistance $R=\qty{40}{\ohm}$. Quelle
intensité $I$ la traverse ?
\end{exo}

\begin{exo}[loi d'Ohm : déduire la résistance]
Un dipôle ohmique soumis à une tension $U=\qty{9}{\volt}$ est parcouru par un courant
$I=\qty{0.5}{\ampere}$. Quelle est sa résistance $R$ ?
\end{exo}

\begin{exo}[brancher un voltmètre]
On veut mesurer la tension aux bornes de la lampe $L$ du circuit ci-dessous. Faut-il brancher le
voltmètre \textbf{en série} ou \textbf{en parallèle} avec $L$ ? Et un ampèremètre, comment se
branche-t-il ?
\begin{center}
\begin{circuitikz}[scale=0.9,transform shape]
  \draw (0,0) to[battery1,l_=$G$] (0,2) -- (2,2) to[lamp,l=$L$] (4,2) -- (4,0) -- (0,0);
\end{circuitikz}
\end{center}
\end{exo}

\begin{exo}[additivité des tensions en série]
Deux résistances $R_1$ et $R_2$ sont en série. On mesure $U_1=\qty{4}{\volt}$ aux bornes de $R_1$ et
$U_2=\qty{6}{\volt}$ aux bornes de $R_2$. Quelle est la tension $U_{AB}$ aux bornes de l'ensemble ?
\begin{center}
\begin{circuitikz}[scale=0.9,transform shape,>=Stealth]
  \draw (0,0) to[short,o-] (0.6,0) to[R,l=$R_1$] (2.6,0) to[R,l=$R_2$] (4.6,0) to[short,-o] (5.2,0);
  \node[above,font=\footnotesize] at (1.6,0.05){$U_1$}; \node[above,font=\footnotesize] at (3.6,0.05){$U_2$};
  \draw[defcol,<->,line width=1pt] (0,0.9) -- (5.2,0.9) node[midway,above,font=\footnotesize,defcol]{$U_{AB}$};
  \node[below,font=\footnotesize] at (0,0){$A$}; \node[below,font=\footnotesize] at (5.2,0){$B$};
\end{circuitikz}
\end{center}
\end{exo}

\end{document}
